\documentclass{article}
% NeurIPS 2026 Education Track: single-blind (author names included) and the
% "education, final" template variant requested in the Call for Educational Resources.
\usepackage[education, final]{neurips_2026}
\usepackage[utf8]{inputenc}
\usepackage[T1]{fontenc}
\usepackage{hyperref}
\usepackage{url}
\usepackage{booktabs}
\usepackage{amsmath}
\usepackage{enumitem}

\title{Breaking GRPO: An Interactive Failure Lab for\\Reinforcement Learning with Verifiable Rewards}

\author{
  Samaksh Gulati \\
  \texttt{samaksh.gulati@dell.com}
  \And
  Rachel Shalom \\
  \texttt{rachel.shalom@dell.com}
}

\begin{document}

\maketitle

\begin{abstract}
Group Relative Policy Optimization (GRPO) became the default algorithm for reinforcement learning with verifiable rewards almost overnight, yet every existing resource teaches only \emph{how} it works. We teach the concept the field actually converged on in 2025: \emph{GRPO's structural failure modes, and the successor design space they produced}. Learners drive three failures --- zero-gradient groups, entropy collapse, and length bias --- in a CPU-only interactive lab framed on KernelBench GPU-kernel tasks, then implement the corresponding DAPO fixes and watch each failure disappear. The result is a failure-to-successor bridge that no current GRPO resource provides.
\end{abstract}

\section{The Concept}

\textbf{Concept taught: the failure modes of GRPO and the 2025 successor design space (DAPO, DisCO, S-GRPO).} This is an emerging concept, not an established one: GRPO was introduced in 2024~\cite{deepseekmath}, scaled by DeepSeek-R1 in 2025~\cite{deepseekr1}, and its documented failure modes directly produced three NeurIPS 2025 papers.

GRPO removes two of the four models in the PPO/RLHF pipeline: the learned reward model is replaced by a deterministic verifier, and the critic is replaced by the mean reward of a group of $G$ sampled responses. The group \emph{is} the baseline, so advantages are $A_i = (R_i - \text{mean}(R)) / (\text{std}(R) + \epsilon)$.

That one substitution is also the source of its fragility, and this is the insight learners take away:

\begin{center}
\emph{GRPO learns only when its samples contain informative contrasts --- and fails when the group or the verifier provides none.}
\end{center}

Three failures follow mechanically from that sentence, at three different structural levels, and each maps one-to-one onto a technique in DAPO~\cite{dapo} (Table~\ref{tab:map}). Learners do not read about the failures; they set the sliders that cause them. No existing resource (HuggingFace's GRPO tutorial, DeepLearning.AI's RL module, community walkthroughs) provides this failure taxonomy or connects it to the post-GRPO literature.

\begin{table}[h]
\centering
\small
\begin{tabular}{@{}llll@{}}
\toprule
\textbf{Failure mode} & \textbf{Level} & \textbf{What the learner triggers} & \textbf{Successor fix} \\
\midrule
Zero-gradient groups & Group & All rollouts agree $\Rightarrow \sigma=0 \Rightarrow A_i=0$ & Dynamic Sampling \\
Entropy collapse & Token & Symmetric clip starves rare good tokens & Clip-Higher \\
Length bias & Sequence & $1/|o_i|$ over-weights short answers & Token-Level Loss \\
Truncation noise & Sequence & Hard length cutoff creates a reward cliff & Overlong Shaping \\
\bottomrule
\end{tabular}
\caption{The failure-to-successor map. Each lab experiment produces a failure that motivates a specific DAPO technique, which learners then implement.}
\label{tab:map}
\end{table}

\section{Leveling and Prerequisite Knowledge}

\textbf{Level: advanced undergraduate to graduate / practitioner.} Prerequisites are probability and expectations, gradients, the softmax function, and how an LLM samples tokens autoregressively. \textbf{No prior reinforcement learning knowledge is assumed}: a short scaffolding section ($\sim$10 of $\sim$55 minutes) establishes the SFT$\to$RL$\to$PPO vocabulary needed to make the failure analysis legible. That scaffolding is deliberately brief and is not the contribution --- the emerging material is the failure lab and the successor methods.

\section{Learning Objectives and Outcomes}

By the end, a learner will be able to:
\begin{enumerate}[nosep,leftmargin=1.4em]
  \item \textbf{Understand} why replacing PPO's critic with a group mean is both GRPO's efficiency win (4 models $\to$ 2) and the origin of its failure modes.
  \item \textbf{Apply} the group-relative advantage by hand and implement a working GRPO update in $\sim$30 lines of NumPy.
  \item \textbf{Analyse} why symmetric ratio clipping is asymmetric in probability space, and why that starves exploration.
  \item \textbf{Evaluate} when zero-gradient groups will dominate, predicting the wasted-compute fraction from success rate and group size $G$.
  \item \textbf{Create} the four DAPO fixes and verify each repairs its corresponding failure, plus design a reward-hacking verifier and a defence against it.
  \item \textbf{Evaluate} whether GRPO, DAPO, DisCO, or S-GRPO suits a given task using a readiness checklist.
\end{enumerate}

\section{Linked Papers}

The concept is in active use at the frontier. GRPO originates in DeepSeekMath~\cite{deepseekmath} and powers DeepSeek-R1~\cite{deepseekr1}. Three NeurIPS 2025 papers build directly on it, and each is motivated by a failure this lab reproduces: \textbf{DAPO}~\cite{dapo} introduces Dynamic Sampling, Clip-Higher, Token-Level Loss and Overlong Reward Shaping, explicitly citing entropy collapse and reward noise; \textbf{DisCO}~\cite{disco} removes the question-difficulty bias in group-relative normalisation via a discriminative objective and constrained optimisation; \textbf{S-GRPO}~\cite{sgrpo} replaces parallel groups with serial exit points and decaying rewards, cutting reasoning length 40--61\%. The verifiable-reward setting we use for framing is KernelBench~\cite{kernelbench} (ICML 2025), and Kevin~\cite{kevin} applies multi-turn RL over those same tasks, independently reporting entropy collapse and reward hacking at scale --- evidence that the toy failures are not artifacts.

\section{Teaching Materials Summary}

\begin{itemize}[nosep,leftmargin=1.4em]
  \item \textbf{Notebook} (\texttt{grpo\_failure\_lab.ipynb}): 65 cells (29 code, 36 markdown), $\sim$55 minutes, pure NumPy/Matplotlib. Three interactive \texttt{ipywidgets} experiments, each with a static fallback so the notebook renders without a kernel. Runs end-to-end on a laptop CPU in under a minute; no GPU, LLM, or API key.
  \item \textbf{Two micro-environments} framed on KernelBench: a \emph{Micro-Bandit} (choose the tile/block configuration for the Square-Matmul problem) isolating group-level failures, and a \emph{Micro-Sequence} (emit a 4-stage tiled-kernel pipeline) isolating token- and sequence-level failures.
  \item \textbf{Figures}: 15 generated plots, including the PPO-vs-GRPO pipeline comparison, the clipping-asymmetry single-step walkthrough, entropy-collapse curves, and the DAPO-vs-GRPO repair comparison.
  \item \textbf{Exercises}: 6 analytical, implementation, and diagnostic exercises with collapsible in-notebook solutions plus a standalone verified solution script.
  \item \textbf{Data}: KernelBench Level-1 task specifications referenced by real problem ID (\texttt{1\_Square\_matrix\_multiplication\_}, \texttt{23\_Softmax}, \texttt{40\_LayerNorm}). Per-kernel latency/speedup values are clearly labelled \emph{illustrative} --- representative of published KernelBench and Kevin findings, not live \texttt{nvcc} runs --- with provenance and GPU reproduction instructions documented alongside the data.
  \item \textbf{Supporting}: a plain-language primer for pre-reading, and this summary. All materials are original and created for this track.
\end{itemize}

\bibliographystyle{plainnat}
\begin{thebibliography}{9}
\small

\bibitem[Shao et~al.(2024)]{deepseekmath}
Z.~Shao, P.~Wang, Q.~Zhu, R.~Xu, J.~Song, M.~Zhang, Y.~Li, Y.~Wu, and D.~Guo.
\newblock DeepSeekMath: Pushing the limits of mathematical reasoning in open language models.
\newblock \emph{arXiv:2402.03300}, 2024.

\bibitem[DeepSeek-AI(2025)]{deepseekr1}
DeepSeek-AI.
\newblock DeepSeek-R1: Incentivizing reasoning capability in LLMs via reinforcement learning.
\newblock \emph{arXiv:2501.12948}, 2025.

\bibitem[Yu et~al.(2025)]{dapo}
Q.~Yu, Z.~Zhang, R.~Zhu, Y.~Yuan, X.~Zuo, Y.~Yue, et~al.
\newblock DAPO: An open-source LLM reinforcement learning system at scale.
\newblock In \emph{NeurIPS}, 2025. arXiv:2503.14476.

\bibitem[Li et~al.(2025)]{disco}
G.~Li, M.~Lin, T.~Galanti, Z.~Tu, and T.~Yang.
\newblock DisCO: Reinforcing large reasoning models with discriminative constrained optimization.
\newblock In \emph{NeurIPS}, 2025. arXiv:2505.12366.

\bibitem[Dai et~al.(2025)]{sgrpo}
M.~Dai, C.~Yang, and Q.~Si.
\newblock S-GRPO: Early exit via reinforcement learning in reasoning models.
\newblock In \emph{NeurIPS}, 2025. arXiv:2505.07686.

\bibitem[Ouyang et~al.(2025)]{kernelbench}
A.~Ouyang, S.~Guo, S.~Arora, A.~L.~Zhang, W.~Hu, C.~R\'e, and A.~Mirhoseini.
\newblock KernelBench: Can LLMs write efficient GPU kernels?
\newblock In \emph{ICML}, 2025. arXiv:2502.10517.

\bibitem[Baronio et~al.(2025)]{kevin}
C.~Baronio, P.~Marsella, B.~Pan, S.~Guo, and S.~Alberti.
\newblock Kevin: Multi-turn RL for generating CUDA kernels.
\newblock \emph{arXiv:2507.11948}, 2025.

\bibitem[Schulman et~al.(2017)]{ppo}
J.~Schulman, F.~Wolski, P.~Dhariwal, A.~Radford, and O.~Klimov.
\newblock Proximal policy optimization algorithms.
\newblock \emph{arXiv:1707.06347}, 2017.

\end{thebibliography}

\end{document}
